\begin{table}[H]
\centering
\caption{Ranking departamental de niveles de remuneraci\'on por trabajador, 2024}
\label{tab:pib_geih_productividad_departamento_remuneracion_niveles}
\scriptsize
\begin{tabular}{lrrrr}
\toprule
Departamento & Rem./trab. 2024 & Puesto & PIB/trab. 2024pr & Puesto \\
\midrule
Bogotá D.C. & 2,95 & 1 & 67,4 & 1 \\
Antioquia & 2,12 & 2 & 47,7 & 5 \\
Caldas & 1,85 & 3 & 36,6 & 12 \\
Meta & 1,83 & 4 & 63,1 & 2 \\
Risaralda & 1,82 & 5 & 38,7 & 10 \\
Santander & 1,80 & 6 & 62,0 & 3 \\
Valle del Cauca & 1,79 & 7 & 54,5 & 4 \\
Cundinamarca & 1,79 & 8 & 40,4 & 8 \\
Quindío & 1,71 & 9 & 33,8 & 13 \\
Atlántico & 1,58 & 10 & 37,3 & 11 \\
Tolima & 1,50 & 11 & 40,0 & 9 \\
Norte de Santander & 1,49 & 12 & 24,3 & 20 \\
Boyacá & 1,49 & 13 & 47,5 & 6 \\
Caquetá & 1,47 & 14 & 24,9 & 18 \\
Huila & 1,36 & 15 & 31,6 & 15 \\
Bolívar & 1,33 & 16 & 45,1 & 7 \\
Cesar & 1,26 & 17 & 32,6 & 14 \\
Chocó & 1,16 & 18 & 28,6 & 16 \\
Magdalena & 1,15 & 19 & 23,9 & 21 \\
Cauca & 1,06 & 20 & 26,6 & 17 \\
Nariño & 1,06 & 21 & 17,8 & 24 \\
Córdoba & 1,04 & 22 & 22,0 & 22 \\
Sucre & 1,01 & 23 & 22,0 & 23 \\
La Guajira & 0,89 & 24 & 24,8 & 19 \\
\bottomrule
\end{tabular}
\caption*{\footnotesize Nota: remuneraci\'on por trabajador en millones de pesos constantes de 2025 al mes; PIB por trabajador en millones de pesos constantes de 2015. La tabla est\'a ordenada por el nivel de remuneraci\'on por trabajador. Fuente: c\'alculos propios con DANE y GEIH.}
\end{table}